\documentclass[11pt,a4paper]{article}
\usepackage[utf8]{inputenc}
\usepackage[margin=1in]{geometry}
\usepackage{graphicx}
\usepackage{hyperref}
\usepackage{xcolor}
\usepackage{enumitem}
\usepackage{titlesec}
\usepackage{mdframed}
\usepackage{booktabs}
\usepackage{amsmath,amssymb}

\hypersetup{
    colorlinks=true,
    linkcolor=blue,
    urlcolor=blue
}

\titleformat{\section}{\large\bfseries\color{blue!80!black}}{}{0pt}{}[\hrule]
\titleformat{\subsection}{\bfseries\color{teal}}{}{0pt}{}

\mdfdefinestyle{pitchstyle}{
    linecolor=blue!80!black,
    backgroundcolor=blue!5,
    leftmargin=10pt,
    rightmargin=10pt,
    innerleftmargin=10pt,
    innerrightmargin=10pt,
    innertopmargin=8pt,
    innerbottommargin=8pt,
    linewidth=1pt
}

\title{\textbf{Conceptual Framework \& Multi-Domain Applications Guide}\\ \large Proposed Cognitive Graph Digital Twin (GDT) Adaptive Cyber-Security Architecture}
\author{\textbf{Anshika Keshari}*\textsuperscript{1}, \textbf{Mallikharjuna Rao K}\textsuperscript{2}, \textbf{Sabarathinam Chockalingam}\textsuperscript{3}, \textbf{Sandeep Pirbhulal}\textsuperscript{4} \\
\small \textsuperscript{1}KIIT Bhubaneswar \& IIIT Naya Raipur, \textsuperscript{2}IIIT Naya Raipur, \textsuperscript{3}IFE Norway, \textsuperscript{4}NR Norway}
\date{\today}

\begin{document}

\maketitle

\tableofcontents
\newpage

\section{Overview of the Proposed Conceptual Architecture}

The proposed \textbf{Cognitive Graph Digital Twin (GDT) Adaptive Cyber-Security Framework} operates across three interconnected pillars:

\begin{enumerate}
    \item \textbf{Physical Cyber-Physical Twin (IT-OT Domain)}: Comprises SCADA control centers, industrial 5G networks, and field Operational Technology (OT) assets.
    \item \textbf{Virtual Graph Digital Twin Layer}: Real-time spatial graph mirror $G(V, E)$ tracking nodal features ($V, \theta, P, Q$) and admittance attributes ($Y_{\text{bus}}$).
    \item \textbf{Cognitive Decision-Making \& Healing Cycle}: Four-stage Observe-Learn-Plan-Act (O-L-P-A) loop executing Multi-Task GraphSAGE detection, Integrated Gradients explainability, and closed-loop self-healing.
\end{enumerate}

\section{Detailed Breakdown of the Three Core Pillars}

\subsection{Physical Cyber-Physical Twin (IT-OT Domain)}
\begin{itemize}
    \item \textbf{IT Layer}: SCADA master stations, cloud management servers, and energy management systems (EMS).
    \item \textbf{Industrial Network Layer}: High-speed communication protocols (5G slices, IEC 61850) streaming sensor telemetry.
    \item \textbf{OT Field Layer}: Physical hardware assets including Remote Terminal Units (RTUs), Phasor Measurement Units (PMUs), busbar substations, and circuit breakers.
\end{itemize}

\subsection{Virtual Graph Digital Twin Layer}
Spatial admittance graph representation $G = (V, E, X, E_{\text{attr}})$ where buses are nodes, transmission lines are edges, and nodal admittance $Y_{\text{bus}} = G_{\text{bus}} + jB_{\text{bus}}$ forms the graph topology. Real-time synchronization allows the virtual graph to dynamically update its edge connections when lines trip ($N-1$ outages).

\subsection{Cognitive Decision-Making \& Adaptive Security Cycle (O-L-P-A)}
\begin{enumerate}
    \item \textbf{Observe \& Orient}: Normalizes raw SCADA sensor streams ($V, \theta, P, Q$) and constructs spatial PyTorch Geometric graph objects.
    \item \textbf{Learn}: Multi-task GraphSAGE GNN backbone executes concatenation-based message passing:
    \begin{equation}
    \mathbf{h}_v^{(k)} = \text{ReLU}\left(\mathbf{W}_k \cdot \left[\mathbf{h}_v^{(k-1)} \,\|\, \mathbf{h}_{\mathcal{N}(v)}^{(k)}\right]\right)
    \end{equation}
    Predicts graph-level attack presence ($98.25\%$ F1) and node anomalies ($97.62\%$ F1) in $0.908$~ms.
    \item \textbf{Plan \& Explain}: Integrated Gradients (IG) computes exact per-sensor attributions ($\sum \text{Attr}_i = \Delta F$).
    \item \textbf{Decide \& Act}: Quarantines corrupted RTU channels and substitutes load baseline forecasts ($\pm 2\%$ noise), reducing state estimation MSE from $0.007228\text{ p.u.}^2$ down to $0.000002\text{ p.u.}^2$ (\textbf{99.97\% error recovery}).
\end{enumerate}

\section{Multi-Domain CPS Applications}

While our paper focuses on \textbf{Smart Power Grids (IEEE 14-bus / 30-bus)}, the Graph Digital Twin framework is domain-agnostic:

\begin{table}[htbp]
  \centering
  \resizebox{\linewidth}{!}{%
  \begin{tabular}{llll}
    \hline
    \textbf{CPS Application} & \textbf{Graph Nodes ($V$)} & \textbf{Graph Edges ($E$)} & \textbf{Telemetry Signals ($X$)} \\
    \hline
    \textbf{Smart Power Grids} & Transmission Substations & Power Lines & Voltage ($V$), Angle ($\theta$), Power ($P, Q$) \\
    \textbf{Water Distribution Networks} & Pipe Junctions \& Pumps & Water Main Pipes & Water Flow Rate, Hydraulic Pressure \\
    \textbf{Industrial Microgrids} & Battery / Solar Inverters & Microgrid Busbars & State-of-Charge (SoC), Frequency, Power \\
    \textbf{Smart EV Charging Grids} & Fast Charging Depots & Substation Feeders & Charging Demand, Transformer Voltage \\
    \hline
  \end{tabular}%
  }
\end{table}

\end{document}
